\subsection{Measurement details}\label{app:measurement}

\paragraph{Prediction criteria.} Unseen KL is computed on test positions
whose preceding pair of boards never occurred at a training position; its
support is recorded per run (about 20,000--26,000 positions). Witness
recovery is $1 - \bar{K}/J$, where $\bar{K}$ is the mean KL in bits on the 64
histories made of an anchor and one \Ncat{} board, and $J = 0.061278\ldots$
bits is the Jensen--Shannon divergence between the two laws, the mean KL of a
predictor that ignores the anchor. The state-intervention panels reset and
set the mode, hold it through neutral boards, substitute a board of the same
category and exchange upper states. Under the equal law, witness recovery is
undefined and is replaced by a check that the two witness predictions agree
within 0.02 TV. r10 uses its own conjunction, including bars derived
from its calibrated history-blind predictor and law TV at most 0.05 under
its original law; the complete study-specific criteria for every study,
including r10's covered-value KL, equal-law TV and witness-separation rules,
are listed in Supplement Tables S4, S8, S10, S12 and S14 (r8, r10,
r11, r12, r13; r9 has no pass criteria).

\paragraph{Census.} The two contexts are produced by executing an \Lcat{} and
an \Hcat{} anchor through the network, so a misread can also reflect a
failure to establish the mode. The signature distance is the summed absolute
difference over the two contexts. From r11 on, the census covers all 24,435
boards in each of four fixed renderings and counts are reported per
rendering; r9 used one rendering per board.

\paragraph{Probes.} In the original task the upper state is probed after one
board from a zero state; in r10 it is probed after recurrent histories. The
2,048 probe-training boards and the 512 evaluation boards per slot are
disjoint from each other but not necessarily unseen during network training.
A positive gain is supported when the lower bound of its paired 95\%
interval (1,000 resamples of the same 512 boards) is above zero and, for
linear probes, both fits converged; an interval wholly below zero supports a
decline. Linear probes use balanced class weights, which can push ordinary
accuracy below the majority floor.

\paragraph{Cutoff discrimination (supplementary).} A separate analysis of
the probability-target records $\to$ prediction networks shows that probes
separate the sums 12 from 13 (150 evaluation pairs) and 23 from 24 (13
pairs) in every seed, while those networks still fail registered prediction
tests. This establishes discrimination on that support, not across all
cutoff boards, and does not diagnose the numerical-encoding failures.

\paragraph{Swap test.} For the networks that learn from records, the
alignment is a mean-centred orthogonal map, fitted on training data only,
from the 16-dimensional raw carrier onto 15 coordinates that
track the exact sums (three per slot) plus one residual coordinate; the
patched state is mapped back through its inverse. It is a fixed, fitted
alignment, not a searched subspace. The trained positive (records+sums
$\to$ prediction) is instead patched by replacing the queried slot's hard
sum answer directly. The r10 swap test uses a fixed panel of 1,024 board
pairs per slot and requires at least 20 pairs per slot on which the network
is accurate for both boards. Donor-follow is the fraction of those accurate
pairs for which the patched prediction is closer in TV to the donor's law
than to the base board's law. Selectivity requires that an independent
decoder finds no change in the other four slots' sums. Four checks are
reported separately: whether the trained positive (at the 20,000-update
endpoint) follows the donor, whether the no-op and shuffled-alignment
negatives do not, whether the swap is selective, and whether the evaluated
networks follow the donor.
